\begin{frame}{도메인 무작위화: 범위를 무작위화하라, 최악을 최소화하라}
  ``시뮬레이션 점수만 보고 고르는 것''의 함정: 실물의 $\mu$ 를
  \emph{정확히} 알 수 없으니까.
  \vskip0.4em
  \kb{도메인 무작위화}는 이 문제를 학습 \emph{중}에 푼다: $\mu$ 를
  매 에피소드마다 $\{0.1, 1, 2, 5\}$ 처럼 \emph{범위}에서 무작위로 뽑아
  훈련 $\to$ ``$\mu = 0.1$ 인 세계''에 과적합할 기회 자체가 사라짐.
  \vskip0.4em
  \small
  \begin{center}
    시뮬레이션 최고점 정책 $(k_p\!=\!8, k_d\!=\!2)$ 의 $\mu$ 별 점수:
    $\;0.1$: $-1163.9$ \;\; $1.0$: $-1225.9$ \;\; $2.0$: $-1225.1$ \;\;
    $5.0$: $-1204.3$
  \end{center}
  \vskip0.5em
  \begin{block}{가장 위험한 환경은 예측하지 못한 곳}
    시뮬레이션(0.1)도, 가정한 실물(5.0)도 아닌
    \textbf{$\mu = 1.0 \sim 2.0$ 이라는 중간 값}.
    $\to$ 어느 $\mu$ 에서도 나쁘지 않은 정책 = \kb{``모든 $\mu$ 에서 가장
    늦게 떨어지는'' ($k_p = 8$ 계열)} --- 이 절의 모든 이야기가 이 한 줄로
    요약됨. 정규화 강도 $\lambda$ 와 정확히 같은 \textbf{트레이드오프}.
  \end{block}
\end{frame}
